\documentclass{ctexart}
\usepackage{amsmath, amssymb}
\usepackage{geometry}
\geometry{a4paper, margin=2cm}

\title{\textbf{第3章 功和能} --- 知识点与公式}
\author{}
\date{}

\begin{document}
\maketitle

\section*{3.1 功}

\subsection*{一、恒力的功}
\[
A = Fs\cos\theta = \vec{F}\cdot\vec{S}
\]

\subsection*{二、变力的功}
元功：
\[
\mathrm{d}A = \vec{F}\cdot\mathrm{d}\vec{r}
\]
从 $a$ 到 $b$ 沿路径 $L$ 的总功：
\[
A = \int_{a(L)}^{b} \vec{F}\cdot\mathrm{d}\vec{r}
\]

直角坐标系：
\[
A = \int_{a(L)}^{b} (F_x\,\mathrm{d}x + F_y\,\mathrm{d}y + F_z\,\mathrm{d}z)
\]

自然坐标系（$|\mathrm{d}\vec{r}| = \mathrm{d}s$）：
\[
A = \int_{a(L)}^{b} F\cos\theta\,\mathrm{d}s
\]

\textbf{说明：}
\begin{itemize}
  \item 功是标量，有正负
  \item 合力的功等于各分力的功的代数和
  \item 功的值一般与路径有关
\end{itemize}

\subsection*{三、功率}
平均功率：
\[
P = \frac{\Delta A}{\Delta t}
\]
瞬时功率：
\[
P = \frac{\mathrm{d}A}{\mathrm{d}t} = \vec{F}\cdot\vec{v} = Fv\cos\theta
\]

\section*{3.2 几种常见力的功}

\subsection*{一、重力的功}
\[
A = \int_{z_1}^{z_2} (-mg)\,\mathrm{d}z = mg(z_1 - z_2)
\]
结论：重力的功只与始末位置的高度差有关，与路径无关。

\subsection*{二、弹性力的功}
弹簧弹性力 $F = -kx$：
\[
A = \int_{x_1}^{x_2} -kx\,\mathrm{d}x = \frac12 kx_1^{2} - \frac12 kx_2^{2}
\]
结论：弹性力的功只与始末位置有关，与路径无关。

\subsection*{三、万有引力的功}
\[
\mathrm{d}A = -G\frac{mM}{r^{2}}\,\mathrm{d}r
\]
\[
A = \int_{r_1}^{r_2} -G\frac{mM}{r^{2}}\,\mathrm{d}r = GmM\left(\frac{1}{r_2} - \frac{1}{r_1}\right)
\]
结论：万有引力的功只与始末位置有关，与路径无关。

\subsection*{四、摩擦力的功}
\[
A = -\mu mg s
\]
结论：摩擦力的功与路径有关。

\section*{3.3 动能定理}

\subsection*{一、质点动能定理}
\[
\mathrm{d}A = \vec{F}\cdot\mathrm{d}\vec{r} = m\vec{v}\cdot\mathrm{d}\vec{v}
\]
\[
A = \frac12 mv_2^{2} - \frac12 mv_1^{2} = E_{k2} - E_{k1}
\]
作用于质点的合力在某一路程中所作的功，等于质点在同一路程始末两状态动能的增量。

说明：
\begin{itemize}
  \item 动能 $E_k = \dfrac12 mv^{2}$ 是标量，是状态量
  \item 功是过程量
  \item 合力作正功时动能增加；合力作负功时动能减少
  \item 动能定理只适用于惯性系
\end{itemize}

\subsection*{二、质点系动能定理}
\[
\sum_i A_i = \sum_i \frac12 m_i v_{i2}^{2} - \sum_i \frac12 m_i v_{i1}^{2}
\]
\[
A_{\text{外}} + A_{\text{内}} = \Delta E_k
\]

\textbf{注意：}内力和为零，但内力功的和不一定为零。功是标量，其和为代数和。成对内力可改变系统的动能（如炸弹爆炸）。

\section*{3.4 势能 机械能守恒定律}

\subsection*{一、保守力}
如果力所做的功与路径无关，只决定于物体的始末相对位置，这样的力称为保守力。
保守力沿闭合路径一周所做的功为零：
\[
\oint_{L} \vec{F}\cdot\mathrm{d}\vec{r} = 0
\]
\begin{itemize}
  \item 保守力：重力、万有引力、弹性力
  \item 非保守力：摩擦力
\end{itemize}

\subsection*{二、势能}
在保守力场中，质点从某点到势能零点，保守力所作功的大小为该点的势能：
\[
E_p = \int_{M}^{M_0} \vec{F}\cdot\mathrm{d}\vec{r}
\]

\textbf{重力势能}（选地面为零点）：
\[
E_p = mgy
\]

\textbf{弹性势能}（选原长为零点）：
\[
E_p = \frac12 kx^{2}
\]

\textbf{万有引力势能}（选无穷远为零点）：
\[
E_p = -G\frac{mM}{r}
\]

注意：
\begin{itemize}
  \item 势能大小由相对位置决定，势能零点选取是任意的
  \item 保守力作正功时，动能增加、势能减少
  \item 保守力作负功时，动能减少、势能增加
  \item 保守力的功 $A = -\Delta E_p$
\end{itemize}

\subsection*{三、势能曲线与保守力}
保守力是势能的负梯度：
\[
\vec{F} = -\nabla E_p = -\left(\frac{\partial E_p}{\partial x}\vec{i} + \frac{\partial E_p}{\partial y}\vec{j} + \frac{\partial E_p}{\partial z}\vec{k}\right)
\]
势能曲线上某点斜率的负值，就是该点对应的保守力。

判断保守力的三维条件（交叉偏导相等）：
\[
\frac{\partial F_x}{\partial y} = \frac{\partial F_y}{\partial x},\quad
\frac{\partial F_y}{\partial z} = \frac{\partial F_z}{\partial y},\quad
\frac{\partial F_z}{\partial x} = \frac{\partial F_x}{\partial z}
\]

\subsection*{四、机械能守恒定律}
对质点系：
\[
A_{\text{外}} + A_{\text{非内}} = \Delta E_k + \Delta E_p = \Delta E
\]
当 $A_{\text{外}} + A_{\text{非内}} = 0$ 时，$\Delta E = 0$，即机械能守恒。

说明：
\begin{itemize}
  \item 守恒条件：外力和非保守内力作功之和为零
  \item 守恒定律是对一个系统而言的
  \item 守恒是对整个过程而言的，不能只考虑始末两状态
\end{itemize}

\section*{3.5 能量守恒定律}
能量不能消失，也不能创造，只能从一种形式转换为另一种形式。对一个封闭系统，各种形式的能量可以互相转换，但总和保持不变。

\begin{itemize}
  \item 能量守恒定律适用于任何变化过程
  \item 功是能量交换或转换的一种度量
  \item 机械能守恒定律是能量守恒定律在机械运动范围内的体现
\end{itemize}

\end{document}
